\DocumentMetadata{
  pdfversion=2.0,
  pdfstandard=ua-2,
  lang=en-US,
  tagging=on,
  tagging-setup={math/setup=mathml-SE,tabsorder=structure}
}
\documentclass[11pt,letterpaper]{article}
\usepackage[margin=0.68in,includefoot]{geometry}
\usepackage{newcomputermodern}
\usepackage{amsmath,amscd,mathtools}
\usepackage{tikz,adjustbox}
\providecommand{\square}{\mdlgwhtsquare}
\usepackage[hidelinks]{hyperref}
\hypersetup{
  pdftitle={Algebra PhD Exam, May 2000},
  pdfauthor={University of Florida Department of Mathematics},
  pdfsubject={Graduate mathematics examination},
  pdfdisplaydoctitle=true,
  bookmarksopen=true
}
\setcounter{secnumdepth}{0}
\setlength{\parindent}{0pt}
\setlength{\parskip}{0.28em}
\setlength{\emergencystretch}{3em}
\setlength{\leftmargini}{2.15em}
\setlength{\leftmarginii}{2.65em}
\setlength{\leftmarginiii}{3em}
\renewcommand{\labelenumi}{\textbf{\theenumi.}}
\pagestyle{empty}
\raggedbottom
\allowdisplaybreaks
\newenvironment{examproblems}[1]{%
  \begin{enumerate}%
  \setcounter{enumi}{#1}%
  \setlength{\topsep}{0.35em}%
  \setlength{\partopsep}{0pt}%
  \setlength{\itemsep}{0.48em}%
  \setlength{\parsep}{0.18em}%
}{\end{enumerate}}
\newenvironment{examsubparts}{%
  \begin{enumerate}%
  \setlength{\topsep}{0.18em}%
  \setlength{\partopsep}{0pt}%
  \setlength{\itemsep}{0.16em}%
  \setlength{\parsep}{0.1em}%
}{\end{enumerate}}
\newenvironment{examinstructions}{%
  \par\begingroup\itshape
}{%
  \par\endgroup\vspace{0.18em}
}
\makeatletter
\renewcommand\section{\@startsection{section}{1}{0pt}%
  {-1.25ex plus -.25ex minus -.1ex}%
  {0.55ex plus .1ex}%
  {\normalfont\large\bfseries}}
\renewcommand{\@maketitle}{%
  \newpage\null\vskip -1em%
  {\footnotesize\bfseries
    \href{https://www.ufl.edu/}{University of Florida}, \href{https://math.ufl.edu/}{Department of Mathematics}\par}%
  \vskip -0.5em%
  \tagstructbegin{tag=Title}%
    {\LARGE\bfseries\href{https://gma.math.ufl.edu/exams/algebra-phd/}{Algebra PhD Exam}, May 2000\par}%
  \tagstructend%
  \vskip 0.25em%
  \tagmcbegin{artifact}%
    \hrule height 1pt%
  \tagmcend%
  \vskip 1em%
}
\makeatother
\title{Algebra PhD Exam, May 2000}
\author{}
\date{}
\begin{document}
\maketitle
\thispagestyle{empty}
\begin{examinstructions}
Please work out seven out of the following eleven exercises. Please do not hand in more than seven.
\end{examinstructions}
\begin{examproblems}{0}
\item
Let~\(F\) denote the free group on the set~\(X\). Prove that \(F/[F,F]\) is the free abelian group on the set~\(X\). (Note: \([F,F]\) denotes the commutator subgroup of~\(F\).)
\item
Recall that a subgroup~\(H\) of a product of groups \(G=\prod_{i\in I}G_i\) is said to be a subdirect product if, for each \(i\in I\), the restriction of the projection \(\pi_i:G\longrightarrow G_i\) to~\(H\) is surjective.

\par
Let \(\mathcal V\) be a class of groups which is closed under the formation of subdirect products and homomorphic images. Prove that \(\mathcal V\) is also closed under formation of subgroups.
\item
Let~\(R\) be a ring with identity, and~\(M\) be a left~\(R\)-module. A projective basis for~\(M\) is a subset \(\{a_i:i\in I\}\) of~\(M\) together with~\(R\)-homomorphisms \(f_i:M\longrightarrow R\) such that the following conditions hold. Prove that~\(M\) is a projective module if and only if it has a projective basis.
\begin{examsubparts}
\item[\textnormal{(i)}]
for each \(a\in M\), \(\left|\{i\in I:f_i(a)\neq 0\}\right|<\infty\), and
\item[\textnormal{(ii)}]
for each \(a\in M\), \(a=\sum_{i\in I}f_i(a)a_i\).
\end{examsubparts}
\item
Suppose that~\(G\) is a finite group and~\(K\) is a field. Prove:
\begin{examsubparts}
\item[\textnormal{(a)}]
If the characteristic of~\(K\) does not divide \(|G|\), then \(K[G]\), the group algebra over~\(K\), is a semisimple left Artinian ring. (You may use that \(K[G]\) is a ring with identity, and the Wedderburn–Artin Theorem in all its glory.)
\item[\textnormal{(b)}]
Show that if the characteristic of~\(K\) does divide the order of~\(G\), then the Jacobson radical of \(K[G]\) is nontrivial.
\end{examsubparts}
\item
Let~\(A\) be a commutative ring with identity, and suppose that~\(M\) is an~\(A\)-module, and~\(I\) is an ideal of~\(A\).
\begin{examsubparts}
\item[\textnormal{(a)}]
Prove that 
\[
(A/I)\otimes M\cong M/IM,
\]
 where \(IM\) is the submodule generated by all elements of the form \(xb\), with \(x\in I\), \(b\in M\).
\item[\textnormal{(b)}]
Prove that the isomorphism in (a) is natural.
\end{examsubparts}
\item
Suppose that~\(A\) is a commutative ring with identity. If every prime ideal of~\(A\) is finitely generated, show that~\(A\) is Noetherian.
\item
State and prove the Hilbert Basis Theorem.
\item
Let~\(K\) be a field.
\begin{examsubparts}
\item[\textnormal{(a)}]
Define discrete valuation on~\(K\), and the notion of a discrete valuation ring.
\item[\textnormal{(b)}]
Suppose that~\(A\) is an integral domain. Prove that it is a discrete valuation ring if and only if it is a principal ideal domain in which 
\[
p\text{ and }q\text{ irreducible}\implies p=uq,
\]
 for some unit \(u\in A\).
\end{examsubparts}
\item
Prove that the lattice of all ideals of a Dedekind domain is distributive. (Hint: Localize!)
\item
Use the notion of formal derivatives to show that, if \(f(T)\) is an irreducible polynomial over the field~\(F\), then it has repeated roots in the splitting field if and only if the characteristic of~\(F\) is \(p>0\), and \(f(T)=g(T^p)\), for some \(g(T)\in F[T]\).
\item
Let~\(p\) be a prime number.
\begin{examsubparts}
\item[\textnormal{(a)}]
Prove that if a subgroup~\(H\) of \(S_p\), the symmetric group on~\(p\) letters, contains a~\(p\)-cycle and a transposition, then \(H=S_p\).
\item[\textnormal{(b)}]
If \(p(T)\) is an irreducible polynomial over \(\mathbb Q\), of degree~\(p\), having exactly two non-real roots, then show that the Galois group of the splitting field of \(p(T)\) is \(S_p\).
\end{examsubparts}
\end{examproblems}
\end{document}
